\scanpage{6}
Avec l'application source donnée par $\operatorname{pr}_2$, l'application but par la loi d'opération de $\mathcal E$ sur $\varepsilon$ (via $\gamma$ et $\mathcal E\longrightarrow\gamma$), la loi de composition étant donnée par
\[
(u,vx)\circ(v,x)=(uv,x)
\]
Considérons un groupe $\gamma'$, un \deleted{\unsure{hom}} $\gamma'$-torseur $\varepsilon'$, et un homomorphisme de torseurs
\[
\varepsilon'\;-\;\varepsilon
\]
associé à un hom de groupes
\[
\gamma'\;-\;\gamma,
\]
on veut décrire le \inserted{$\pi$-}groupoïde induit sur $\varepsilon'$ comme ens. d'objets par le $\pi$-groupoïde $\mathcal E(\varepsilon)$, grâce à l'application précédente. On trouve qu'il est can$^{\mathrm t}$ isom au $\pi$-groupoïde $\mathcal E'(\varepsilon')$, où $\mathcal E'$
\[
(11)\qquad0-\pi\longrightarrow\mathcal E'\longrightarrow\gamma'\longrightarrow0
\]
est l'extension, image inverse de (8) par $\gamma'\;-\;\gamma$.

Revenant alors à la description de $S_0\mathcal T_{0,I}$, en termes de l'inclusion (3), on voit qu'on a
\[
(12)\qquad S_0\mathcal T_{0,I}\deleted{\unsure{=}}\simeq\mathcal E_I(\omega(I))
\]
\margin{Dans les \inserted{\unsure{conn}} $\mathbb Z^I$-groupoïdes}
où $\mathcal E_I$ est l'extension
\[
(13)\qquad0-\mathbb Z^I\longrightarrow\mathcal E_I\longrightarrow\{\pm1\}\longrightarrow0
\]
\mspage{5}
